% Original: PDF strana 11, gornji slajd.
\slajdid{03-s021}
\begin{frame}[label=03-s021]{Lanac kodera prioriteta sa 16 ulaza}
\begin{columns}[c,onlytextwidth]
\begin{column}{.56\textwidth}\centering
% element: 03-s021:e01
\begingroup
\def\DEenable{1}
% \KPSimbol{ime}{x}{y}; pozivaoci koriste milimetarsku mrežu, font 9 pt.
\newcommand{\KPSimbol}[3]{%
\begin{scope}[shift={(#2,#3)}]
\draw[DEvod] (0,0) rectangle (22,12);
\ifnum\DEenable=1
 \def\DEgsheight{9}
 \node at (11,6.9) {KP};
 \node[anchor=north,inner sep=.7pt] at (11,12) {EI};
 \draw[DEvod] (11,12)--(11,12.6);
 \coordinate (#1-ei) at (11,12.6);
 \node[anchor=south,inner sep=.7pt] at (11,0) {EO};
 \draw[DEvod] (11,0)--(11,-.6);
 \coordinate (#1-eo) at (11,-.6);
 \foreach \n/\y in {1/5.4,0/2.2}{
  \node[anchor=east,inner sep=.7pt] at (22,\y) {A\n};
  \draw[DEvod] (22,\y)--(24,\y);
  \coordinate (#1-a\n) at (24,\y);
 }
\else
 \def\DEgsheight{6.5}
 \node at (11,10.5) {KP};
 \foreach \n/\x in {1/8,0/15}{
  \node[anchor=south,inner sep=.7pt] at (\x,0) {A\n};
  \draw[DEvod] (\x,0)--(\x,-1.5);
  \coordinate (#1-a\n) at (\x,-1.5);
 }
\fi
\foreach \n/\y in {3/9.5,2/7,1/4.5,0/2}{
 \node[anchor=west,inner sep=.7pt] at (0,\y) {I\n};
 \draw[DEvod] (-2,\y)--(0,\y);
 \coordinate (#1-i\n) at (-2,\y);
}
\node[anchor=east,inner sep=.7pt] at (22,\DEgsheight) {GS};
\draw[DEvod] (22,\DEgsheight)--(24,\DEgsheight);
\coordinate (#1-gs) at (24,\DEgsheight);
\end{scope}}

\begin{tikzpicture}[circuit logic US,x=.93mm,y=1mm,font=\fontsize{9}{11}\selectfont,
 zbir/.style={or gate US,draw,logic gate inputs=nnnn,minimum width=7mm,minimum height=8mm}]
\foreach \k/\y in {3/40.5,2/27,1/13.5,0/0}{
 \KPSimbol{p\k}{0}{\y}
 \foreach \n in {3,2,1,0}{
  \pgfmathtruncatemacro{\ii}{4*\k+\n}
  \draw[DEvod] (p\k-i\n)--++(-1.5,0) node[left,inner sep=.7pt] {I\ii};
 }
}
\foreach \vis/\niz in {3/2,2/1,1/0}
 \draw[DEvod] (p\vis-eo)--(p\niz-ei);
\draw[DEvod] (p3-ei)--++(0,1.5) node[above,inner sep=.7pt] {1};
\KPSimbol{f}{48}{27}
\draw[DEvod] (f-ei)--++(0,1.5) node[above,inner sep=.7pt] {1};
\foreach \k/\x in {3/44,2/42,1/42,0/44}
 \draw[DEvod] (p\k-gs)--(\x,0 |- p\k-gs)|-(f-i\k);
\draw[DEvod] (f-gs)--++(1.5,0) node[right,inner sep=.7pt] {GS};
\draw[DEvod] (f-a1)--++(1.5,0) node[right,inner sep=.7pt] {A3};
\draw[DEvod] (f-a0)--++(1.5,0) node[right,inner sep=.7pt] {A2};
\draw[DEvod] (f-eo)--++(0,-2)--++(14,0) node[right,inner sep=.7pt] {EO};
\node[zbir] (a1) at (56,14) {};
\node[zbir] (a0) at (56,4) {};
\foreach \k/\pin/\xx in {3/1/29,2/2/31,1/3/33,0/4/35}
 \draw[DEvod] (p\k-a1)--(\xx,0 |- p\k-a1)|-(a1.input \pin);
\foreach \k/\pin/\xx in {3/1/30,2/2/32,1/3/34,0/4/36}
 \draw[DEvod] (p\k-a0)--(\xx,0 |- p\k-a0)|-(a0.input \pin);
\draw[DEvod] (a1.output)--++(2.5,0) node[right,inner sep=1pt] {A1};
\draw[DEvod] (a0.output)--++(2.5,0) node[right,inner sep=1pt] {A0};
\end{tikzpicture}
\endgroup

\end{column}
\begin{column}{.40\textwidth}
% element: 03-s021:e02
Lanac \(EI/EO\) propušta dozvolu ka koderima nižeg prioriteta dok se ne naiđe na aktivan ulaz. \(A_1,A_0\) iz osnovnih kodera spajaju se logičkim ILI kolima.
\par\vspace{2mm}
% element: 03-s021:e03
Ako neaktivan koder postavlja \(A_1,A_0\) u visoku impedansu, izlazi se mogu spojiti bez tih ILI kola.
\par\vspace{2mm}
% element: 03-s021:e04
\Oznaka{Konvencija simbola}
Ulazi su obično s jedne, a izlazi s druge strane pravougaonog simbola; svi signali moraju biti označeni.
\end{column}
\end{columns}
\end{frame}
